\section{引言}\label{sec:intro}

在协同巡检、按需无线通信中继与密集编队重构等自主任务中，多机器人系统往往同时受制于两类方向相反、相互冲突的标量耦合安全约束：一方面，成对机间避碰约束要求相邻机器人之间保持安全排斥距离；另一方面，全局网络连通性保持约束（以图拉普拉斯代数连通度标量下界表征）将全网机器人耦合在一起，限制编队过度分散以防通信链路割裂\cite{yang2010decentralized,sabattini2013decentralized,ong2021connectivity}。控制障碍函数（Control Barrier Function, CBF）\cite{ames2017cbf,breeden2022sampled}通过将安全集前向不变性条件转化为关于控制输入的标量仿射不等式，为安全关键控制提供了高计算效率的二次规划（Quadratic Program, QP）滤波框架。针对包含全局耦合行与成对接触行的标量安全约束，现有研究已发展出基于连续时间无违约分布式梯度流\cite{tan2025violationfree}以及基于零和配额分解的邻域近端分配（Neighborhood Proximal Allocation, NPA）方法，使得多机器人系统在每个采样时刻仅凭有限轮一跳邻居通信即可获得严格满足全部耦合约束的安全控制输入。

然而，单拍安全滤波器的瞬时可行性并不蕴含长时程多机导航任务的完成性。当多台机器人在狭窄空间相向交会或需要在高密度阵列内部交换位置时，指向各自目标的标称吸引力往往被接触避碰约束与全局连通约束的法向反力同时抵消，使系统永久停滞在远离目标的对称平衡构型（即几何死锁，Geometric Deadlock）\cite{reis2021,grover2023deadlock,tan2024equilibria,mestres2026filters}。从物理与凸优化机制上看，传统 CBF 最小修正滤波器之所以在密集交会中频繁陷入死锁，根源在于其“逐拍单调降低原始位置误差”的短视结构。记全网机器人当前位置 $\bp$ 与固定标签目标 $\bg$ 的总欧氏位置误差为 $W(\bp)=\frac12\sum_{i=1}^N\norm{\bp_i-\bg_i}_2^2$。将指向目标的负梯度标称速度 $-k_g\nabla_{\bp}W=-k_g(\bp-\bg)$ 投影到包含零速度解的安全输入集上，必然强制要求 $\dot W(\bp(t))\le 0$\cite{allibhoy2024safegradient}。在内部节点迎面互换而外围邻居已处于各自目标位置（$\bp_j=\bg_j$）的密集编队中，这一单调性约束会产生双重锁死：其一，待互换节点的径向目标引力与直接接触避碰法向反力共线抵消，缺乏侧向破对称驱动；其二，外围静止邻居一旦暂时离开原位腾出几何通道就会导致总误差 $W(\bp)$ 短暂上升，且在静态配额平摊下其局部可行域在接触边界上直接坍缩为零\cite{karagoz2014navigation}。

为缓解安全滤波器引发的伪平衡点与死锁问题，现有文献主要探索了启发式切向环流与学习型安全控制两条技术路线。在解析控制方面，Reis 等\cite{reis2021}与 Tan 等\cite{tan2024equilibria}分析了单智能体及多智能体 CBF-QP 的非目标平衡点结构与鞍点不稳定性；Gon\c{c}alves 等\cite{goncalves2023circulation}与 Dai 等\cite{dai2024compatible}分别引入成对切向环流项与兼容控制李雅普诺夫—障碍函数（CLF-CBF）以打破单对接触下的共线停滞。然而，当内部互换节点被多层已达标的静止邻居紧密包裹时，无差别的常值环流会在多条接触边上引发旋转方向冲突，且在存在严格约束局部极小（而非鞍点）的对称接触构型下，任何依赖局部负曲率或要求 $W$ 单调不增的解析反馈均在数学上无法越过能量壁垒。在学习控制方面，原始—对偶拉格朗日约束多智能体强化学习（MAPPO-Lagrangian）\cite{gu2023safe}、预测与自适应安全滤波强化学习\cite{wabersich2021,liu2026saferl}、神经图控制障碍函数（GCBF+）\cite{zhang2025gcbf}及离散图策略近端优化（DGPPO）\cite{zhang2025dgppo}利用图神经网络（Graph Neural Network, GNN）学习多机协同与避碰策略。但现有多智能体安全强化学习方法仍面临三项悬而未决的结构性挑战：
\begin{enumerate}
 \item \textbf{全局与局部冲突耦合约束下的随时零违约保证缺失：} 神经障碍函数与软惩罚策略梯度仅能近似处理依赖局部相对坐标的成对避碰约束，无法在训练探索与有限通信轮数截断执行下严格保证全局标量连通约束 $\lambda_2(\bL(\bp))\ge\lr$ 与成对避碰约束在采样区间内的连续时间硬安全；
 \item \textbf{学习标称层与分布式动态配额分配层的脱节：} 现有安全护盾多采用单机独立投影或静态责任平摊（$B=0$），当外围机器人已到达自身目标（$\be_j=\boldsymbol 0$）时，静态平摊护盾会将其向外退让通道锁死，迫使学习策略在标称层同时拟合全部外围节点的退让轨迹，在多机密集接触下信用分配效率极低；
 \item \textbf{密集目标构型处的零阶排斥残差悬停：} 消息传递 GNN 输出的成对相对方向排斥项属于不随位置误差 $\be_k\to\boldsymbol 0$ 自动衰减的零阶偏置；在接触净空仅为毫米级的密集晶格终点处，持续非零的接触对偶乘子使冲突门控保持常开，导致机器人在目标微球外围发生持续稳态悬停与微幅抖振。
\end{enumerate}

针对上述挑战，本文提出包络等变图策略与邻域近端分配（Envelope Equivariant Graph Policy with Neighborhood Proximal Allocation, EGP-NPA）协同控制方法，将复杂的多机协同脱困问题解耦为“一跳等变图网络残差标称 + 闭邻域块动态分摊护盾 + 一跳闭邻域终点完成包络”三个紧密咬合的功能模块。本文的主要贡献归纳如下：
\begin{enumerate}
 \item \textbf{严格非目标局部极小的精确有理数二阶 KKT 证书与脱困物理机制（定理\ref{thm:seven-kkt}）：} 针对全员可动的七机器人双三角对称接触构型，利用有理数域代数闭合信道核与精确分数算术（$\mathbb Q$）构造原始—对偶 KKT 见证，严格证明拉格朗日 Hessian 在完整临界方向锥上一致正定（$\boldsymbol v^{\top}\boldsymbol H_L\boldsymbol v\ge\frac{19}{515}\norm{\boldsymbol v}_2^2$），并给出半径 $1/200$ 球面上的显式二次能量壁垒 $W(\bar\bp+\boldsymbol\delta)-W(\bar\bp)\ge\frac1{100}\norm{\boldsymbol\delta}_2^2>0$（在初态与目标均严格位于安全集内部的微扰开球上依然成立）及满足全程 $\lambda_2(\bL)\ge 1792/625>5/2$ 的解析可达路径。该证书从代数上排除了依靠鞍点负曲率或单调下降控制实现一般脱困的可能性，并定量揭示两项核心脱困物理机制：待互换节点的侧向破对称标称驱动与外围静止邻居经由动态配额让渡的协同退让。
 \item \textbf{手性旋转与置换等变的一跳图神经网络残差标称架构（命题\ref{prop:equivariance}）：} 在局部目标正交标架与有向边正交标架上构造融合几何投影与滤波器内部代数状态（单机重拟合碰撞影子价格 $\boldsymbol\mu_{i,k-1}$、谱配额与边记录）的两层参数共享消息传递 GNN，严格证明标称控制场对平面刚体变换群 $\mathrm{SO}(2)\ltimes\mathbb R^2$ 与节点置换群 $S_N$ 的等变性；同时从代数与实测两方面证明：保留特殊正交群 $\mathrm{SO}(2)$ 的确定性手性旋转偏置是打破镜像共线死锁的必要结构条件，若强制施加全正交群 $\mathrm{O}(2)$ 镜像反射等变，则在对称迎头交会构型上切向旋转输出必恒等于零（$\omega_{i,k}^\theta\equiv 0$）并退化回短视死锁。
 \item \textbf{闭邻域动态份额解冻闭式解与一跳闭邻域终点完成包络理论（命题\ref{prop:safety}--\ref{prop:envelope}）：} 证明在任意探索策略与有限轮跨规模部署下闭环连续轨迹的随时硬安全不变性（命题\ref{prop:safety}）；导出接触边界上闭邻域块二次规划将移动节点的法向切入分量 $\delta_n>0$ 转化为碰撞份额转移 $|t_{ij}^\star|=\frac{\eta\ds}{\eta+4\ds^2}\delta_n>0$ 及静止邻居向外退让速度 $-\boldsymbol n_{ij}^{\top}\bu_j^\star=\frac{\eta}{2(\eta+4\ds^2)}\delta_n>0$ 的显式闭式解（命题\ref{prop:unfreeze}），揭示静态配额平摊（$B=0$）导致人为冻结以及在策略损失中惩罚法向滤波修正量反而抹平退让动力 $\delta_n$ 的机理；提出仅依赖一跳目标误差的闭邻域完成包络 $\rho_{i,k}$，证明其在保持协同阶段残差幅值的同时将终点零阶排斥偏置一阶压制为 $\mathcal O(\norm{\be_k}_2)$，恢复目标构型的零速度平衡点（命题\ref{prop:envelope}）。
 \item \textbf{小规模短轮数训练、跨规模零样本迁移与二阶物理底盘验证：} 基于小规模短轮数（$N\in\{6,\dots,10\},B=2$）多智能体近端策略优化（MAPPO）训练的单一策略检查点，在训练规模（$N=8$，$64$ 组）与跨规模零样本迁移（$N\in\{12,20,25\},B=8$，$32$ 组）共 $96$ 组密集晶格基准及 $48$ 组独立盲测实例上实现 $100\%$ 达标、零步长减半与零有理数审计违约。进一步结合固定权重单变量干预消融与接入在线无除法谱估计器的 \texttt{MuJoCo} 二阶全向底盘闭环实验，验证了各模块的因果贡献并刻画了密集格点下物理安全裕度内缩与通道可通行性的权衡边界（定理\ref{thm:passability-tradeoff}）。
\end{enumerate}

\textbf{符号说明：} $\mathbb R^n$、$\mathbb Q^n$ 与 $\mathbb S_+^n$ 分别表示 $n$ 维实向量空间、有理向量空间与 $n\times n$ 实对称半正定矩阵锥；$\I_n$ 为 $n$ 阶单位矩阵，$\boldsymbol 1_N\in\mathbb R^N$ 为全 $1$ 向量，$\boldsymbol\varepsilon_i\in\mathbb R^N$ 为第 $i$ 个标准基向量；$\norm{\cdot}_2$ 表示向量欧氏范数或矩阵谱范数；$\col(\boldsymbol x_1,\dots,\boldsymbol x_m)=[\boldsymbol x_1^{\top},\dots,\boldsymbol x_m^{\top}]^{\top}$ 表示列堆叠算子；$\otimes$ 表示 Kronecker 积；$(x)_+=\max(x,0)$；$\clip(x,a,b)=\min(\max(x,a),b)$；$\bJ=\begin{bmatrix}0&-1\\1&0\end{bmatrix}$ 表示平面逆时针 $90^\circ$ 正交旋转矩阵，满足 $\bJ^{\top}=-\bJ$ 与 $\bJ^2=-\I_2$；$\mathbb I(\cdot)$ 为示性函数。
